\section{Introduction}
\label{sec:intro}

A single biological sequence model should ideally answer qualitatively different questions
through one interface, and the interface can be as simple as text: present the sequence
\texttt{GTATTCTCCCCTAACAAGAC...} with the question \emph{``Is this DNA sequence a
promoter?''} and the candidates \{yes, no\}, and a model returns $P(\text{yes})$; swap only
the question and candidate set --- \{All Alpha, All Beta, Alpha and Beta, ...\} for fold
class, five value-anchored ordinal levels for mutation fitness, per-label yes/no for
subcellular localization --- and the same model, with no task-specific output heads and no
free-text generation, answers all of them. \emph{Typed-decision} interfaces implement
exactly this design. They are being adopted in
practice, by commercial systems \citep{almeida2026jev} and by open reimplementations
\citep{gliclass2025,nanojev2026,semif2025}, on the promise of one small model, dynamic
candidate sets, and zero-adaptation-cost new tasks. Yet two basic questions have never been
answered under controlled conditions. \textbf{Q1:} does a shared typed-decision scorer match
or beat conventional per-task output heads at matched compute? \textbf{Q2:} does the
``dynamic candidates / zero adaptation'' selling point translate into a measurable
performance or cost advantage?

Answering these questions runs into three obstacles. First, biological sequence evaluation
is fragmented: GUE \citep{zhou2024dnabert}, Genomic Benchmarks
\citep{gresova2022genomic}, TAPE \citep{rao2019tape}, ProteinGym \citep{notin2023proteingym}
and friends each define their own task formats, splits and metrics, so decision architectures
cannot be compared in one protocol. Second, leakage auditing is largely absent: several
upstream ``train'' pools on public hubs \emph{contain the blind-test members of prior
studies}, so a random re-split silently fabricates a new test set. Third --- the
methodological spine of this paper --- \emph{evaluation variance in this regime has never
been quantified}. In our small-budget multi-task setting, two runs of the same configuration
and seed differ by $\pm0.05$--$0.1$ on task-macro metrics, enough to reverse the winner of
an architecture comparison (\cref{sec:mainresults}). Any single-run architectural conclusion
is therefore unreliable, including our own early one.

\paragraph{Approach.}
We build a benchmark first and a measurement protocol second, then compare architectures.
BioDecisionBench converts ten source families (290 tasks; DNA, protein, RNA, and double
sequences) into one typed-decision record format with a leak-proof admission protocol
(\cref{sec:benchmark}). On a 37-task matched slice we train the two comparable arms --- a
shared typed-decision scorer (\arch{1}) and per-task heads on the same encoder (\arch{2})
--- with family-balanced interleaving, stabilized optimization, per-family dev early
stopping, and three seeds; the frozen Qwen3-0.6B full-candidate likelihood (\arch{3}) serves
as an off-the-shelf reference (\cref{sec:framework}). \Cref{sec:results} reports the
seed-averaged comparison and three operational tests of \arch{1}'s claimed value, plus
training-dynamics and mechanism ablations.

\paragraph{Contributions.}
\begin{itemize}
\item \textbf{BioDecisionBench} (\cref{sec:benchmark}): a 290-task, three-modality,
four-primitive biological decision benchmark with a leak-proof admission protocol --- exact
cross-split sequence isolation, conflict-group quarantine (1{,}984 groups, not
majority-voted), verbatim preservation of legacy blind tests, and label mappings validated
on full pool$\cap$test overlaps. Post-isolation within-task leakage is zero.
\item \textbf{A seed-averaged evaluation protocol with a documented failure of single runs}
(\cref{sec:seedavg,sec:mainresults}): same-configuration reruns differ by up to
$7\times$ on score Spearman (0.107 vs 0.015) and reverse the comparison outcome; we make
$\ge3$ seeds, per-task seed means, SEM + run-std reporting, and dev early stopping the
minimum bar for this class of comparison.
\item \textbf{The first matched-compute three-way comparison of decision architectures on
biology} (\cref{sec:mainresults,sec:value}): per-task heads $\ge$ shared scorer (Noul
significantly better, Choice/Score at parity, and more stable across seeds); both trained
arms beat the frozen candidate-likelihood baseline. The shared scorer's advantage is
\emph{qualitative} --- it can serve arbitrary new tasks with zero adaptation, but its
zero-shot quality equals the untrained base, cross-family transfer is weak and
classification-only, and ${\approx}32$ labels let a fresh head catch its zero-shot level.
\item \textbf{Training-dynamics and mechanism diagnostics} (\cref{sec:dynamics,sec:ablations}):
naive budget scaling destabilizes joint training below the untrained baseline; stabilization
recovers but confirms a plateau; per-family checkpoint selection is a zero-inference-cost
improvement (family-optimal steps span $26\times$); two reusable controlled ablations
(noul candidate phrasing across three stable checkpoints; contrastive decoding for frozen
VLM candidate likelihood).
\end{itemize}

\paragraph{Results preview.}
Seed-averaged, per-task heads reach Noul AUROC 0.629 against the shared scorer's 0.478
(SEMs disjoint) --- while a single earlier run had suggested the opposite ordering. The
shared scorer's zero-shot interface, its headline feature, is caught by a fresh 32-label
linear head on 7 of 11 held-out tasks.

\Cref{fig:hero} summarizes the benchmark and the headline comparison.

\begin{figure*}[t]
\centering
\includegraphics[width=0.90\textwidth]{fig1_hero.pdf}
\caption{\textbf{BioDecisionBench and its headline comparison.}
\emph{Left:} 290 leak-audited tasks over 10 source families and 4 decision
primitives; the admission protocol leaves zero within-task train/test sequence
overlap (1{,}984 conflicting-label groups quarantined rather than majority-voted).
\emph{Right:} seed-averaged (3 seeds), matched-compute comparison per primitive
(mean$\pm$SEM over tasks). Per-task heads (\#2) beat the shared typed-decision
scorer (\#1) on noul (disjoint SEMs) and tie on choice/score; both trained arms
exceed the frozen candidate-likelihood baseline (\#3, not compute-matched).
\emph{Inset:} the methodological warning --- two single runs of the same
configuration give score Spearman 0.107 vs 0.015 for \#1; the 3-seed mean is 0.013.}
\label{fig:hero}
\end{figure*}
